\documentclass{beamer}

\usepackage{amsmath, amssymb}

\title{Inference for Categorical Time Series with Missing Observations}
\author{Victor Minig}
\date{27.04.2026}

\begin{document}

\frame{\titlepage}

%------------------------------------------------
\begin{frame}{Motivation}
\begin{itemize}
    \item Many real-world time series are \textbf{categorical}:
    \begin{itemize}
        \item Credit ratings (AAA, AA, ...)
        \item Weather states (sunny, rainy, ...)
        \item Survey responses (agree, neutral, disagree)
    \end{itemize}
    
    \vspace{0.3cm}
    
    \item Classical time series methods assume \textbf{numeric structure}
    
    \vspace{0.3cm}
    
    \item $\Rightarrow$ Need new tools for:
    \begin{itemize}
        \item Computing proper tests for marginal features 
        \item Measuring temporal dependence without bias 
    \end{itemize}
\end{itemize}
\end{frame}

%------------------------------------------------
\begin{frame}{Categorical Time Series}
\begin{block}{Definition}
A categorical time series $(X_t)$ takes values in a finite set
\[
S = \{s_0, \dots, s_m\}.
\]
\end{block}

\vspace{0.3cm}

\textbf{Two types of categorical data:}
\begin{itemize}
    \item \textbf{Nominal:} no natural ordering
    \item \textbf{Ordinal:} inherent ordering
\end{itemize}

\vspace{0.3cm}

\textbf{Key implication:}
\begin{itemize}
    \item Available statistical tools depend on structure of $S$
\end{itemize}
\end{frame}
%------------------------------------------------
\begin{frame}{Marginal Distributions}
\textbf{Central object: marginal probabilities}

\vspace{0.3cm}

\textbf{Nominal data:}
\[
p_i = P(X_t = s_i)
\]
\begin{itemize}
    \item Probability Mass Function (PMF)
\end{itemize}

\vspace{0.5cm}

\textbf{Ordinal data:}
\[
f_i = P(X_t \leq s_i)
\]
\begin{itemize}
    \item Cumulative Distribution Function (CDF)
\end{itemize}

\vspace{0.3cm}

\textbf{Key difference:}
\begin{itemize}
    \item Ordinal $\Rightarrow$ uses ordering
    \item Nominal $\Rightarrow$ no cumulative structure
\end{itemize}
\end{frame}

%------------------------------------------------
\begin{frame}{Location Measures}
\textbf{Nominal data:}
\begin{itemize}
    \item No ordering $\Rightarrow$ no mean/median
    \item Use \textbf{mode}:
\[
\text{Mod}(X_t) = \arg\max_{s_i} P(X_t = s_i)
\]
\end{itemize}

\vspace{0.5cm}

\textbf{Ordinal data:}
\begin{itemize}
    \item Ordering allows additional measures
    \item \textbf{Median:}
\[
\text{Med}(X_t) = \min \{ s_i : F(s_i) \geq 1/2 \}
\]
\end{itemize}

\vspace{0.3cm}

\textbf{Interpretation:}
\begin{itemize}
    \item Mode = most frequent category
    \item Median = central position in ordering
\end{itemize}
\end{frame}

%------------------------------------------------
\begin{frame}{Dispersion Measures}
\textbf{Different notions of variability}
$\Rightarrow$ Same for minimal, different for maximal variability
\vspace{0.5cm}

\textbf{Ordinal: Index of Ordinal Variation (IOV)}
\[
\text{IOV}(\boldsymbol{f}) = \frac{4}{m} \sum_{i=0}^{m-1} f_i (1 - f_i)
\]

\vspace{0.5cm}

\textbf{Nominal: Index of Qualitative Variation (IQV)}
\[
\text{IQV}(\boldsymbol{p}) = \frac{m+1}{m} \left(1 - \sum_{i=0}^{m} p_i^2 \right)
\]

\vspace{0.5cm}

\begin{itemize}
    \item Both scaled to $[0,1]$
\end{itemize}
\end{frame}


%------------------------------------------------
\begin{frame}{Stationarity}
\textbf{Why do we need stationarity?}

\begin{itemize}
    \item Stable probabilistic structure over time
    \item Enables estimation and asymptotics
\end{itemize}

\vspace{0.5cm}

\begin{block}{Strict Stationarity}
\[
F_{X_t, \dots, X_{t+j}} = F_{X_{t+h}, \dots, X_{t+h+j}}
\]
for all $t, j, h$.
\end{block}

\vspace{0.3cm}

\begin{itemize}
    \item Distribution invariant under time shifts
\end{itemize}
\end{frame}




%------------------------------------------------
\begin{frame}{Temporal Dependence}
\textbf{Goal: capture dependence across time}

\vspace{0.3cm}

\textbf{Problem:}
\begin{itemize}
    \item No meaningful arithmetic operations
    \item Classical ACF not applicable
\end{itemize}

\vspace{0.5cm}

\textbf{Solution:}
\begin{itemize}
    \item Work with joint probabilities
\end{itemize}

\vspace{0.3cm}

\begin{block}{Bivariate Probabilities}
\[
p_{ij}(h) = P(X_t = s_i, X_{t+h} = s_j), \quad  f_{ij}(h) = P(X_t \leq s_i, X_{t+h} \leq s_j)
\]
\end{block}

\begin{itemize}
    \item Independence $\Rightarrow p_{ij}(h) = p_i p_j \quad \text{and} \quad  f_{ij}(h) = f_i f_j $
\end{itemize}
\end{frame}

%------------------------------------------------
\begin{frame}{Cohen's $\kappa$}
\textbf{Idea: measure deviation from independence}

\vspace{0.5cm}

\[
\kappa(h)_{\mathrm{nom}} =
\frac{\sum_{j=0}^{m} (p_{jj}(h) - p_j^2)}
{1 - \sum_{i=0}^{m} p_i^2} \quad
\kappa(h)_{\mathrm{ord}} =
\frac{\sum_{j=0}^{m} (f_{jj}(h) - f_j^2)}
{1 - \sum_{i=0}^{m} f_i(1 - f_i)}
\]

\vspace{0.5cm}

\begin{itemize}
    \item $\kappa = 1$: perfect persistence
    \item $\kappa = 0$: independence
    \item $\kappa < 0$: tendency to switch states
\end{itemize}

\vspace{0.5cm}

\textbf{Interpretation:}
\begin{itemize}
    \item Probability of staying in same category
    \item Relative to independence benchmark
\end{itemize}
\end{frame}



%------------------------------------------------
\begin{frame}{Estimation of Marginal Probabilities}
\textbf{Idea: use of indicator functions}

\vspace{0.3cm}

\[
Y_{t,i} = \mathbb{I}\{X_t = s_i\} \quad \text{and} \quad Z_{t,i} = \mathbb{I}\{X_t \leq s_i\}
\]
Then estimators are given by:
\begin{itemize}
    \item $\hat{p}_i = \frac{1}{n} \sum_{t=1}^n Y_{t,i}$
    \item $\hat{f}_i = \frac{1}{n} \sum_{t=1}^n Z_{t,i}$
\end{itemize}

Intuition:\\
$\Rightarrow$ we just count relative frequencies\\
From here we mainly look at ordinal data, but the nominal case is in many ways analogous.
\end{frame}
%------------------------------------------------
\begin{frame}{Mixing Assumptions}

\begin{block}{Idea}
The Central Limit Theorem (CLT) provides asymptotic normality and enables statistical inference.
\end{block}

\begin{alertblock}{Problem}
Classical CLT results assume independence, but in practice data may exhibit serial dependence.
\end{alertblock}

\begin{block}{Solution}
To handle dependence, we impose \textbf{mixing conditions}:
\begin{itemize}
    \item $\alpha$-mixing quantifies the decay of dependence between past and future observations
    \item Exponential decay of the mixing coefficients allows a CLT for dependent data
\end{itemize}
\end{block}
\end{frame}
%------------------------------------------------
\begin{frame}{CLT under serial dependence}
    \begin{block}{Result}
    \[\sqrt{n}(\hat{f} - f) \Rightarrow \mathcal{N}(0, \boldsymbol{\Sigma})\]
    where the elements of $\boldsymbol{\Sigma}$ are given by:
    \[\sigma_{ij} = \underbrace{\left(f_{\min\{i,j\}}- f_i f_j\right)}_{\text{lag}(0)} + \sum_{h=1}^{\infty} \underbrace{\left(f_{ij}(h) + f_{ji}(h) - 2f_if_j\right)}_{\text{lag}(h)}.\]
    Note that $lag(h) \rightarrow 0$ as $h \rightarrow \infty$.
    \end{block}
\end{frame}
%------------------------------------------------
\begin{frame}{CLT's for marginal estimators}
    \begin{block}{Idea}
        IOV is function of the marginal probabilities \\
        $\Rightarrow$ we can apply the \textbf{delta method}  to obtain their asymptotic distributions 
    \end{block}
    \begin{block}{Results}
        \begin{itemize}
            \item Asymptotic variance 
            \[\sqrt{n}Var(\hat{IOV}) \overset{asy}{=} \frac{16}{m^2}\sum_{i,j = 0}^{m-1} (1 - 2f_i)(1 - 2f_j) \sigma_{ij}\]
            \item Bias to order $n^{-1}$:
            \[n Bias(\hat{IOV}) = -\frac{4}{m} \sum_{i = 0}^{m-1}  \sigma_{ii}\]
        \end{itemize}
    \end{block}
\end{frame}
%------------------------------------------------
\begin{frame}{Missingness - Relevance}
    
\end{frame}
%------------------------------------------------
\begin{frame}{Amplitude Modulation}
    \begin{block}{Idea}
        Missingness can be modeled as a binary process $(O_t)$, that \textit{modulates} the \textit{amplitude} of the actual process: 
        \[O_t = 1 \text{ if } X_t \text{ is observed and } O_t = 0 \text{ if } X_t \text{ is missing}\] 
        The observed process is then given by $(O_t X_t)$.
    \end{block}
    \begin{block}{Assumptions}
        \begin{itemize}
            \item $(O_t)$ is stationary with $P(O_t = 1) = \pi > 0$ and joint moments $E[O_t O_{t+h}] = \pi(h)$
            \item $(O_t)$ is $\alpha$-mixing with exponential decay
        \end{itemize}
    \end{block}    
\end{frame}

%------------------------------------------------
\begin{frame}{Missingness Mechanisms}
    Originally developed for cross-sectional data, but also relevant for time series:
    \begin{itemize}  
        \item \textbf{MCAR:} Missing Completely At Random
        \[P(O_t = 1| X_t, X_{t-1}, \ldots)\ = P(O_t = 1)\]
        \item \textbf{MAR:} Missing At Random
        \[P(O_t = 1| X_t, X_{t-1}, \ldots)\ = P(O_t = 1| X_{t-1}, \ldots)\]
        \item \textbf{MNAR:} Missing Not at Random
        \[P(O_t = 1| X_t, X_{t-1}, \ldots)\ \neq P(O_t = 1| X_{t-1}, \ldots)\]
    \end{itemize}
    From here on, unless stated otherwise , we assume MCAR.
\end{frame}

%------------------------------------------------
\begin{frame}{Naive vs corrected estimators}
    \begin{block}{Naive estimator}
        \[\hat{\boldsymbol{\mathfrak{f}}} = \frac{1}{n} \sum_{t=1}^n O_t \boldsymbol{Z}_t\]
        $\Rightarrow$ biased and inconsistent due to missingness
    \end{block}
    \begin{block}{Corrected estimator}
        \[\hat{\boldsymbol{f}}^* = \frac{1}{n} \frac{\sum_{t=1}^n O_t \boldsymbol{Z}_t}{\sum_{t=1}^n O_t}\] 
    \end{block}
\end{frame}
%------------------------------------------------
\begin{frame}{Asymptotic properties of the corrected estimator}
    \begin{block}{Result}
        \[\sqrt{n}(\hat{\boldsymbol{f}}^* - \boldsymbol{f}) \Rightarrow \mathcal{N}(0, \boldsymbol{\Sigma}^*)\]
        where the elements of $\boldsymbol{\Sigma}^*$ are given by:
        \[\sigma^*_{ij} = \frac{1}{\pi} (f_{\mathrm{min}\{i,j\}}- f_i f_j) + \frac{1}{\pi^2} \sum_{h = 1}^{\infty} \pi(h)(f_{ij}(h) + f_{ji}(h) - 2f_if_j)\]
    \end{block}

\end{frame}


\end{document}